\section{Introduction}\label{sec:intro}

A user who asks a model provider to delete their data receives, at best, an assurance
that an unlearning routine was run. Whether the routine removed what it was meant to
remove is a statistical question, and VOUCH \citep{dang2026vouch} gave it a statistical
answer. Before fine-tuning, the provider plants \emph{paired ghost canaries}: pairs of
exchangeable records of which a fair coin admits one into training and, later, into
the forget set, while the other is never trained on. Under exact unlearning the deployed
model cannot depend on the coins, so the sign of the score difference between the two
twins is a fair coin, exactly and at any sample size. Betting e-processes built on these
signs give an anytime-valid $(\eps,\alpha)$-forgetting certificate, a revocation alarm,
and a confidence sequence on the residual advantage, with no shadow models and no
retrained reference.

The certificate VOUCH issues has a precise scope: it bounds the realised advantage of a
declared class of scores on the committed cohort of canaries, under a provider that runs
the same pipeline on canaries and on everything else. That scope leaves three questions
open, and each has practical weight.

\emph{First, do canaries stand in for the data users delete?} VOUCH's canaries are
templated records with random secret strings. Users delete biographies, messages and
facts. A certificate on random strings says nothing about natural-language facts unless
something connects the two, and the connection has to survive a provider who might,
deliberately or not, treat the two differently.

\emph{Second, what does a strategic provider need in order to pass the audit without
unlearning?} The honest-but-verifiable threat model of VOUCH assumes the provider does
not special-case the cohort. The assumption is reasonable for an audit run with the
provider's cooperation, but it is an assumption, and the original study measured that
its own canary templates are easy to recognise. How much can a provider gain by
recognising canaries, and what would it take to deny it that gain?

\emph{Third, does an audit of the trained form see knowledge that survives in other
forms?} Gradient-based unlearning has been reported to suppress the exact string it was
given while the fact it encodes stays reachable through a paraphrase
\citep{li2026beliefs}, and a deletion request may name one record of a fact the model
learned from several. A random-string secret has no paraphrase, so
VOUCH's canaries cannot probe this failure. Semantic canaries can, provided the exact
null survives the move to natural language.

\paragraph{Contributions.} This paper extends VOUCH along these three lines, with new
theory and new experiments on Pythia models \citep{biderman2023pythia}.

\begin{itemize}
\item \textbf{Semantic twins keep the exact null.} A twin pair that states two
incompatible facts about one fictitious entity keeps VOUCH's exact conditional null for
every score that compares the two facts under the same prompt, including belief scores
over paraphrased questions that never appear in training, and their versions calibrated
by the base model (Proposition~\ref{prop:semantic}). This realises the natural-language
canary that the original paper named as its next step.

\item \textbf{A ladder from the cohort to organic records.} Under a bounded-influence
assumption, the cohort certificate transfers to the canary generator's population with a
McDiarmid term (Theorem~\ref{thm:population}); and from canaries to organic records with
a term equal to twice the provider's ability to tell the two apart from their contents
(Theorem~\ref{thm:detect}). The second result prices one deviation from the
honest-provider assumption, content-based special-casing: whatever rule a provider uses to
treat canaries differently, the gap it opens is bounded by that rule's distinguishing
advantage (Corollary~\ref{cor:ladder}).

\item \textbf{Detectability and its price, measured.} VOUCH's random-string records and
semantic twins built from their own templates are detected with advantage one by a
character $n$-gram classifier. Twins drawn from the organic record generator itself, with
one slot varied, are not detected by perplexity or $n$-gram detectors beyond the noise of
the detectors. But when such records appear once, as organic records do, Pythia-410M
retains no detectable trace of them, so the audit becomes valid but powerless: the term
the ladder pays for detectability is traded for power (Sections~\ref{sec:exp-detect}
and~\ref{sec:exp-matched}).

\item \textbf{The literal form and the fact under partial unlearning.} When unlearning
targets only the literal statement, the advantage on paraphrased questions is retained
slightly longer than the literal one in all three experiments we ran (two model sizes, two
designs), by $0.07$ to $0.12$ in like-for-like retention; the direction is consistent but
no test reaches significance at eight seeds, so we do not claim to have demonstrated the
squeezing effect at this scale. What does change verdicts is an incomplete forget set:
when the fact was also trained in a form that the forget set does not contain, a union
audit over literal and calibrated paraphrase scores revokes up to $38\%$ of runs that a
literal-only audit certifies, and it costs nothing in validity when the forget set is
complete (Section~\ref{sec:exp-squeeze}).

\item \textbf{Comparison with membership-inference audits.} On the same runs, a
membership-inference test of canaries against an organic holdout raises a false alarm on
every retrained model, because the test detects the difference between canaries and
organic records rather than membership; against matched ghost twins it is valid but
misses partial unlearning that VOUCH's paired, sequential test detects
(Section~\ref{sec:exp-baselines}).

\item \textbf{When the cohort certificate arrives.} Under a uniformly random reveal order
the betting edge of VOUCH's without-replacement process is deterministic
(Lemma~\ref{lem:edge}). A Freedman bound on the wealth process then gives a tail bound on
the certification time that holds before refutation and shows where the
finite-population speed-up comes from (Theorem~\ref{thm:growth}); it is checked
numerically and is conservative by a factor of two to four.
\end{itemize}

All code, results and the scripts that turn results into the tables of this paper are
released with the VOUCH repository.

\paragraph{What is not claimed.} The organic-record result is conditional on the
distinguishing advantage, which data can bound only from below, and on a bounded-influence
assumption that cannot be checked; we report the strongest detectors we could build, not
a proof that none better exists. Squeezing is not demonstrated. All facts in the
experiments are fictitious, the data come from one domain, and the models have at most
one billion parameters and are fine-tuned with LoRA. These limits are discussed in
Section~\ref{sec:discussion}.

\paragraph{Organisation.} Section~\ref{sec:background} recalls VOUCH and fixes notation.
Section~\ref{sec:related} places the work in the literature.
Sections~\ref{sec:semantic}--\ref{sec:time} develop the theory: semantic twins, the
ladder to organic records, and certification time. Section~\ref{sec:experiments} reports
the experiments and Section~\ref{sec:discussion} discusses limitations.
